\section{序論}
11月14日朝の観測窓では、GoogleのGemini 4シリーズが引き続き中心的な話題となった。高速・低コストを志向するHydrogenの展開に加え、Antigravityで推論モデルと実装向けモデルを切り替える運用が注目された。モデル単体の性能だけでなく、エージェントの実行環境でどのように組み合わせるかが焦点になっている。

同時に、Grok Botの会議参加、Anthropicの安全性評価、気象データセットやMCP拡張など、製品と研究基盤の両面で動きが見られた。フロンティアAIの開発凍結訴訟や既存モデルの品質に関する話題には未確認情報も含まれる。本号はGrok収集の9件を記載順に整理し、独立した一次資料の再検証は行わない。